\chapter*{Abstract}
\addcontentsline{toc}{chapter}{Summary}

% Abstracts are usually around 100–300 words
The Hadley cell (HC) is the meridional large scale circulation governing atmospheric flow in the tropics. Its ascending branch is associated with precipitation and the inter-tropical convergence zone (ITZC), while its descending branch coincides with very dry regions.

With climate change the Hadley circulation is expected to widen and weaken modulating precipitation patterns. Previous studies have shown that this can be mainly attributed to the increase in vertical stability of the atmosphere.
Although CO2 is uniformly distributed in the atmosphere, in our experiments we have decomposed the CO2 forcing into a tropospheric and a stratospheric component using a control tropopause height climatology.

We found the response to be approximately linear with the stratospheric CO2 being responsible for a cooling in the stratosphere and the tropospheric CO2 showing the typical tropospheric warming signal. These temperature anomaly patterns also change the structure of the mean meridional circulation. Unlike general expectations, the HC strength is increasing with tropospheric CO2 while stratospheric CO2 causes a HC weakening.

The angular momentum budget does not provide a satisfactory explanation for the changes in the Hadley circulation. The two limits of angular momentum conservation and eddy driven regimes were not able to explain this behavior, mostly because the flow is never actually close to one of these limits and there is no theoretical solution for the intermediate regime. Instead, we found the energetic framework of the atmosphere to be more insightful. In this view, the HC acts as the primary mechanism by which energy is redistributed within the tropics, thereby reducing the top-of-atmosphere radiative energy imbalance. Its ability to transport energy poleward is controlled by the maximum mass stream function (MSF), the HC strength, as well as by the gross moist stability (GMS) of the atmosphere.  Furthermore, the required energy transport by the HC can be changed by the net energy input into the atmosphere (forcing), the feedback response due to an increase in surface temperature, and the energy partitioning between the mean circulation and the eddies. We find that only by considering a combination of these factors, we can explain the contrasting HC responses.
% We can see that the GMS is the biggest factor to weaken the atmospheric circulation in both cases while the spatial structure of the radiative forcing and the feedbacks determine whether the circulation is weakening or strengthening. 

This study contributes to a more comprehensive understanding of HC response to different patterns of forcing and provides a valuable test of existing theories of atmospheric circulation change.
